\section{Reflection Questions}
\begin{reflection}
\small
\begin{enumerate}[itemsep=0pt, parsep=0pt, topsep=0pt]
  \item When you experience intense negative emotions, do you tend to look inward at your own unmet needs \citep{rizzolatti2004}, or outward to blame someone else?
  \item Which ``pseudo-needs'' (e.g., ``I need you to listen,'' ``I need you to change'') do you catch yourself using most often \citep{streubel2021}?
  \item Think of a recent conflict with your child. What specific strategy were you fighting over, and what was your actual underlying need \citep{lobue2022}?
  \item What fears arise when you consider expressing your deep, vulnerable needs (like a need for security) rather than using demands or judgments \citep{christou2025}?
  \item How did your own caregivers handle their unmet needs? Did they express them clearly, or did they rely on guilt and pseudo-needs \citep{ruba2020}?
  \item In what situations is it hardest for you to separate the behavior (the strategy) from the universal human need driving it \citep{aghaziarati2023}?
  \item How might your relationship change if you responded to your child's anger by asking yourself, ``What are they needing right now?'' \citep{izard2001}
  \item What is one small step you can take today to ensure that your own ``deficit needs'' (like sleep or basic support) are adequately met \citep{herman2025}?
  \item How does practicing self-empathy regarding your own unmet needs impact your ability to be patient with your child's unmet needs \citep{hoffman2000}?
\end{enumerate}
\end{reflection}
